\section{Two-way information comparison and the joint depth--accuracy envelope}
\label{sec:frontier}
Choose two admissible parameters differing only in $b_J$ by $\delta>0$;
$0<\delta<b_+-b_-$ permits such an interior pair. Give both the zero initial
state, an allowed subexperiment of the uniform initial-state class. Write
$q_k=(e_k,0)$, $v_k=(0,e_k)$ and $m_0=2J+1$.

\begin{lemma}[Two-way propagation]
\label{lem:two-way}
For $t\ge0$,
\begin{align}
 |g_\beta(t)-g_\gamma(t)|
 &\le\delta e^{\Lambda t}\Lambda^{4J+2}t^{4J+3}/(4J+3)!,\notag\\
 |h_\beta(t)-h_\gamma(t)|
 &\le\delta e^{\Lambda t}\Lambda^{4J+2}t^{4J+4}/(4J+4)!.
 \label{eq:two-way}
\end{align}
\end{lemma}
\begin{proof}
A generator word from $v_0$ to $q_J$ or from $v_J$ to $q_0$ needs one
velocity-to-position step and two steps for each displacement along an edge.
Hence both boundary entries vanish below degree $m_0$. Damping loops cannot
shorten the path. Each propagator is bounded by
$e^{\Lambda t}(\Lambda t)^{m_0}/m_0!$.
Duhamel's formula for the rank-one difference $A_\beta-A_\gamma=\pm\delta v_Jq_J^*$
is the convolution of these two propagators. The beta integral
$\int_0^t(t-s)^{m_0}s^{m_0}ds=(m_0!)^2t^{2m_0+1}/(2m_0+1)!$
proves the first bound, and integration proves the second.
\end{proof}

For any law on $[0,T]$, \eqref{eq:two-way} gives
$D_\chi\le\delta^2[e^{\Lambda T}\Lambda^{4J+2}T^{4J+4}/(4J+4)!]^2$ for this
pair. Together with \eqref{eq:separation}, this shows, for a fixed positive
admissible $\delta$, that the negative log worst-case $D_\omega$ separation
has order $J\log J$. In reset zero-state windows with bounded force $U$, the
same formula times $nU^2/(2\sigma^2)$ bounds the Gaussian KL divergence.
That is a reset-window testing comparison, \emph{not} a bound on all the
history information in the unreset compact-law experiment. The block
construction explains why the latter can reach a larger depth.

\begin{theorem}[All-time, all-policy information bound]
\label{thm:global}
Let
\[
 \vartheta=\min\{\log2,\lambda_* /(8e\Lambda)\},\qquad
 D_0=(8e\Lambda)^{-1}+C_*/\lambda_*.
\]
For every parameter-independent chronological policy with complete force
bounded by $U$ and $n$ Gaussian boundary readouts at selected finite times,
\begin{equation}
 \KL(P_\beta^{\pi,n},P_\gamma^{\pi,n})
 \le\frac{nU^2\delta^2D_0^4}{2\sigma^2}e^{-4\vartheta(2J+1)}=:K_n.
 \label{eq:kl}
\end{equation}
Every estimator has maximum two-point probability of $d_J$ error at least
$\delta/2$ bounded below by $\frac12(1-\sqrt{K_n/2})$, or zero if negative.
\end{theorem}
\begin{proof}
Split the $L^1$ norm of either off-boundary propagator at
$S=m_0/(8e\Lambda)$. Word vanishing and $m_0!\ge(m_0/e)^{m_0}$ bound the
short part by
$[m_0/(8e\Lambda)]4^{-m_0}\le(8e\Lambda)^{-1}2^{-m_0}$.
The stable tail is at most $(C_*/\lambda_*)e^{-\lambda_*S}$.
Both propagators therefore have $L^1$ norm at most $D_0e^{-\vartheta m_0}$.
Convolving them in Duhamel's formula yields
$\|g_\beta-g_\gamma\|_1\le\delta D_0^2e^{-2\vartheta m_0}$.
At a common observed history, the two candidate inputs are the same bounded
function, so their conditional mean difference at any selected time is at
most $U$ times this norm. The action kernels cancel in the chronological
KL chain rule; each Gaussian readout contributes the squared mean difference
averaged under the true past, divided by $2\sigma^2$.
This proves \eqref{eq:kl}, without conditioning on the terminal design.
Testing the two disjoint strict error balls and using Pinsker's inequality
proves the error bound.
\end{proof}

For fixed accuracy, \eqref{eq:kl} tends to zero when $J_n\gg\log n$.
Together with Theorem~\ref{thm:block-posterior} it proves the order conclusion
of Theorem~\ref{thm:main} within the original bounded-probe, retained-state
experiment with an admissible fast clock. The constructive upper bound is
uniform over the block-committed class; the lower bound covers the larger
fully adaptive bounded-force class.

For the joint scale $J_n=\lfloor r\log n\rfloor$, radius $n^{-s}$, take a
pair separated by $4n^{-s}$ (admissible eventually). Then
$K_n\le Cn^{1-2s-8\vartheta r}$, so uniformly vanishing error at that radius
is excluded when $8\vartheta r+2s>1$.
The constructive sufficient region is $C_b r+p_b s<1$ for bounded probes,
and $C_\infty r+8s<1$ for the finite-mean law. These are explicit inner and
outer regions, not a common exact exponent curve. They imply the two-sided
\emph{order} relation that a logarithmic sample budget controls depth plus
log inverse accuracy, with different box-dependent constants.

Since inter-readout time is at least $\tau_g>0$, a deadline $t$ permits at
most $\lfloor t/\tau_g\rfloor$ readouts. The stopped chronological chain rule
gives the same lower bound with that maximum count. The deterministic upper
clock in the bounded-probe construction makes its attainable depth order
$\log t$ as well. Feedback can change actual information and constants; no
claim of a universally optimal controller or matching leading constants is made.
